% ECONOMICS_PROBLEM: {"id": "C81-368", "title": "LLM-based economic measurement", "jel_code": "C81", "source_id": "OP-0427"}
% !TeX program = xelatex
\documentclass[11pt,a4paper]{article}
\usepackage[margin=23mm]{geometry}
\usepackage{fontspec}
\setmainfont{Latin Modern Roman}
\usepackage{amsmath,amssymb,mathtools}
\usepackage{xcolor,enumitem,booktabs,longtable,array,xurl,hyperref}
\definecolor{muted}{HTML}{53636C}
\hypersetup{colorlinks=true,urlcolor=blue,linkcolor=blue}
\setlength{\parindent}{0pt}
\setlength{\parskip}{5pt}
\newcommand{\R}{\mathbb R}
\newcommand{\E}{\mathbb E}
\newcommand{\N}{\mathbb N}
\newcommand{\ind}{\mathbf 1}
\newcommand{\Var}{\operatorname{Var}}
\newcommand{\Cov}{\operatorname{Cov}}
\newcommand{\supp}{\operatorname{supp}}
\DeclareMathOperator*{\argmin}{arg\,min}
\DeclareMathOperator*{\argmax}{arg\,max}
\newcommand{\entrymeta}[3]{{\sffamily\small\color{muted}\textbf{\texttt{#1}}\quad|\quad #2\par #3\par}}
\newcommand{\Problemref}[1]{\texttt{#1}}
\newcommand{\JELref}[2]{\texttt{#2} (JEL \texttt{#1})}
\begin{document}
\section*{LLM-based economic measurement}
\textbf{Problem ID:} \texttt{C81-368}\quad \textbf{JEL:} \texttt{C81}\par
% BEGIN ECONOMICS STATEMENT
\label{problem:OP-0427}
{\sffamily\small\color{muted}Original subject group: Econometrics, causal inference and measurement\par}
\label{definitions:problem:30007709}
\textbf{Audit classification:} Background; proposed extension. Original2024 and registryv4 December2025 dates are distinguished, and the cited version updated. Fixed-population validation is already analyzed; the new refresh optimization needs a drift law and a width target to avoid a trivial interval solution.
\subsection*{Definitions and precise research target}
Consider a defined population of economic documents with bounded outcome \(Y\in[0,1]\), bounded covariates \(X\), and a human-audited latent measure \(L\in[0,1]\), such as a stated sentiment or contract feature. Model version \(v\) produces automated measure \(\widehat L_v\). The target parameter \(\beta\) is the coefficient on \(L\) in a specified population linear projection of \(Y\) on \(L,X\). Assume the projection is uniquely defined, with residual variance of \(L\) given the covariate basis bounded below by a stated positive constant. A probability validation sample records \(L\), with known positive inclusion probabilities; the remaining documents record \(\widehat L_v,Y,X\). Versions and population composition may change between training, validation and deployment. For a declared class \(\mathcal D\) of shifts, constrain both the change in the measurement-error distribution conditional on \(L,X,Y\) and changes in the deployment law of \(L,X,Y\), using stated total-variation tolerances or representative refreshed validation. Determine the minimum validation design and refresh frequency that permit confidence intervals for deployment \(\beta\) with coverage at least \(1-\alpha\), uniformly over \(\mathcal D\), for stated \(\alpha\in(0,1)\) and audit budget. Distinguish measurement error from training-data leakage and require temporal or private-data checks appropriate to the target documents. Existing validation corrections solve fixed-population cases; unrestricted opaque model drift cannot be corrected from unlabeled text alone. The joint updating, leakage and shift specification here is proposed rather than an explicit unresolved theorem in the cited paper.

\textbf{Definitions, assumptions and mathematical qualifications.} Fix an adjudication procedure \(L=f^*(\text{document})\), rather than claim that an arbitrary human label is error-free truth. Let \(b(X)\) be a bounded finite covariate basis including an intercept, \(R_L=L-\Pi_{b(X)}L\), and \(\beta=E[R_LY]/E[R_L^2]\), with \(E[R_L^2]\geq\kappa>0\). Known validation probabilities satisfy \(\pi_i\geq\pi_{\min}>0\) and refer to a stated probability sampling design. A deployment time \(t\), model version, prompt and decoding rule are recorded. The drift class must bound both changes in \(\widehat L\mid L,X,Y\) and changes in the deployment law of \((L,X,Y)\), or collect representative fresh validation; an error-kernel bound alone does not transport a regression coefficient. Given a width target \(w>0\), require uniform coverage and expected interval length at most \(w\) when asking for minimum validation budget.
\subsection*{Variants and assumption changes}
\begin{enumerate}
\item \textbf{Sourced fixed-population benchmark.} Combine a probability validation sample with automated labels using an appropriate correction; the paper already provides validation-based methods.
\item \textbf{Editorial deployment monitoring.} Allow recorded model updates and bounded joint-population drift over stated refresh intervals; optimize audit timing/budget for a declared coverage-and-width target, including adjudication disagreement sensitivity.
\end{enumerate}
\subsection*{References and source audit}
\textbf{Checked source.} arXiv:2412.07031v4, registry revision 5 December 2025, Sections 3.3, 4.3--4.4 and 6. The paper provides leakage guidance and validation corrections. Full HTML and arXiv registry checked. Original preprint dates to 2024; retrieved HTML has a later 24 August 2026 manuscript dateline, so registry and manuscript dates are distinguished. \url{https://arxiv.org/html/2412.07031v4}.
\begin{enumerate}\raggedright
\item\hypertarget{definitions:source:30007709:1}{} Jens Ludwig; Sendhil Mullainathan; Ashesh Rambachan. 2025. \emph{Large Language Models: An Applied Econometric Framework}. arXiv:2412.07031v4, registry revision 5 December 2025, Sections 3.3, 4.3--4.4 and 6.
\url{https://arxiv.org/html/2412.07031v4}.
DOI: \url{https://doi.org/10.48550/arXiv.2412.07031}.
\end{enumerate}

\subsection*{Common conventions: Source mathematical conventions}

These conventions apply to each entry unless explicitly replaced.

\textbf{Models and identification.} A primitive model \(m\) specifies all probability laws, feasible choices, information, equilibrium and measurement rules used by its target. The admissible class \(\mathcal M\) is fixed independently of outcomes. If \(P_O(m)\) is its observable probability law and \(\tau(m)\) the target, its identified set at observable law \(P\) is
\[
 \mathcal I_\tau(P)=\{\tau(m):m\in\mathcal M,\ P_O(m)=P\}.
\]
Point identification means that this set is a singleton. An empty set signals incompatibility of data and assumptions. Coordinate bounds need not describe the joint set. Bounds are sharp only if every retained target is attainable or arbitrarily approachable by compatible models. Finite bounds require explicit restrictions on unobserved quantities; unrestricted confounding, hidden wealth, extrapolation or arbitrary equilibrium selection can yield uninformative sets.

\textbf{Causal interventions.} A potential outcome \(Y_i(p)\) is the outcome under a complete policy regime \(p\), including timing, financing, information and market spillovers. The target population and units are fixed before treatment. With interference, use the complete assignment vector or a justified exposure mapping; individual no-interference notation alone cannot identify an economy-wide intervention. A difference of mean potential outcomes does not require the cross-world joint distribution. Conditioning on post-treatment migration, survival, enrollment or employment changes the estimand and needs separate justification.

\textbf{Statistical statements.} An estimator, confidence set, forecast or test specifies a sampling law, model class and loss/coverage criterion. Uniform \((1-\alpha)\)-coverage means \(\inf_{m\in\mathcal M}\Pr_m\{\tau(m)\in C_n\}\geq1-\alpha\), or an explicitly asymptotic version. The whole parameter space trivially has coverage, so informativeness requires a stated width, risk or power objective. A forecasting improvement is relative to a named benchmark, information set and evaluation population.

\textbf{Equilibrium and optimization.} An equilibrium comprises feasible plans, optimality, beliefs where needed, budget constraints and all market/resource clearing. The primitives or a stated benchmark define each correspondence \(\mathcal E(p;m)\). Nonemptiness is assumed or proved, never inferred just from compact policy sets. A selected-equilibrium target uses a declared selection rule; a robust target quantifies over all permitted equilibria. Use suprema or infima unless attainment is established. An \(\varepsilon\)-optimal policy is feasible and within \(\varepsilon\) of the finite optimum value. Report an empty feasible set separately.

\textbf{Welfare and accounting.} Normative utility, interpersonal normalization, social weights, numeraire, discounting and the use of public receipts are primitives. Behavioral utility need not coincide with the welfare criterion. Household welfare includes ownership income and rents once; adding profits again to compensating variation generally double-counts them. A monetary equivalent is defined by an explicit transfer and price convention, with monotonicity and range conditions. Average effects, mechanism contributions, transfers and welfare losses are different targets. Infinite sums require convergence and dynamic feasible plans require terminal or no-Ponzi conditions.

\textbf{Source versions.} A working-paper date, revision date and journal publication year are distinguished. A DOI for a journal version is not automatically the DOI of an earlier manuscript. Locators refer to the stated version and printed pages where specified. A metadata-only check does not verify a claim about a full-text conclusion. Every entry's source subsection narrows the provenance claim accordingly.
% END ECONOMICS STATEMENT
\end{document}
